\documentclass{article}
\usepackage{amsmath,amssymb,booktabs,graphicx}
\usepackage[margin=1in]{geometry}

\begin{document}

\title{A One-Site Vector Representation of dry-MARTINI DOPC for Hybrid Upside Membrane Simulations}
\author{}
\date{}
\maketitle

\begin{abstract}
We describe a one-particle coarse-grained lipid (CGL) representation for
implicit-water dry-MARTINI DOPC membranes coupled to the Upside protein model.
Each DOPC molecule is represented by a single particle carrying a dynamic unit
orientation vector.  Conservative CGL interactions are obtained by explicit
projection of dry-MARTINI constituent bead interactions and stored in HDF5
spline tables.  The production tables are generated from isolated DOPC
bonded-conformer ensembles and side-chain rotamer libraries; no bilayer
potential of mean force, iterative Boltzmann inversion, force matching,
thickness restraint, capping, or additional orientation potential is used.  A
fluctuation--dissipation-consistent generalized Langevin equation (GLE) is used
only as a particle-level transport model for CGL diffusion.  The selected
transport parameters give a CGL-only lateral diffusion coefficient of
$0.257~\mu\mathrm{m}^2\,\mathrm{s}^{-1}$ under the Upside-core
$40~\mathrm{ps}$ per step clock, close to the target
$0.25~\mu\mathrm{m}^2\,\mathrm{s}^{-1}$ obtained after applying the
56-fold DOPC coarse-graining factor.  The same force-field tables and transport
parameters produce stable CGL-only bilayers and replicated 1RKL and 1AFO
protein--bilayer hybrid simulations.
\end{abstract}

\section{Model Definition}

\subsection{One-Particle CGL State}

Each dry-MARTINI DOPC molecule is mapped to one CGL particle with position
$\mathbf{R}$ and a dynamic unit vector $\hat{\mathbf{n}}$.  The position is the
geometric center of the 14 dry-MARTINI DOPC beads.  The vector points from the
NC3 head bead toward the midpoint of the terminal C5A and C5B tail beads.
Thus CGL is not an isotropic bead; it is a single translational particle with
an orientational state.

The current implementation does not introduce a hidden marker particle.  The
orientation vector is a non-particle state variable with an associated angular
momentum.  Spline derivatives with respect to angular coordinates are converted
to torques on $\hat{\mathbf{n}}$, and the vector is propagated by rotational
dynamics.  The transverse inertia and orientation length used to define this
rotational state are derived from the isolated DOPC conformer ensemble:
\begin{align}
  \ell &=
  \left\langle
  \left|(\mathbf{x}_{\mathrm{tail}}-\mathbf{x}_{\mathrm{COM}})
  \cdot\hat{\mathbf{n}}\right|
  \right\rangle,\\
  I_\perp &=
  \left\langle
  \sum_i m_i
  \left(
  |\mathbf{x}_i-\mathbf{x}_{\mathrm{COM}}|^2
  -[(\mathbf{x}_i-\mathbf{x}_{\mathrm{COM}})\cdot\hat{\mathbf{n}}]^2
  \right)
  \right\rangle .
\end{align}
For the installed two-conformer DOPC table, the HDF5 metadata records
$\ell=9.406~\text{\AA}$ and
$I_\perp=2.260\times10^4~\mathrm{g\,mol^{-1}\,\text{\AA}^2}$.  These values
are not adjustable fitting parameters; they are consequences of the mapped
dry-MARTINI molecular geometry.

\subsection{Resolved dry-MARTINI Energy}

All conservative CGL tables are generated from explicit dry-MARTINI bead
interactions.  For bead pair $i,j$ separated by distance $d$,
\begin{equation}
  u_{ij}(d) =
  4\epsilon_{ij}
  \left[
  \left(\frac{\sigma_{ij}}{d}\right)^{12}
  -\left(\frac{\sigma_{ij}}{d}\right)^6
  \right]
  +\frac{k_e q_i q_j}{d},
  \qquad d<r_c ,
\end{equation}
and $u_{ij}=0$ for $d\ge r_c$, with $r_c=1.2~\mathrm{nm}$.  Bead types,
charges, Lennard-Jones parameters, masses, and bonded terms are read from the
dry-MARTINI force-field files.  A numerical guard
\begin{equation}
  d_{\mathrm{eval}}=\max(d,10^{-6}~\mathrm{nm})
\end{equation}
is used only to avoid division by zero at exact bead coincidence.

Source dry-MARTINI parameters remain in native units.  Conversion to Upside
runtime units is explicit:
\begin{equation}
  1~E_{\mathrm{up}}=2.914952774272~\mathrm{kJ\,mol^{-1}},
  \qquad
  1~\mathrm{nm}=10~\text{\AA}.
\end{equation}

\section{Force-Field Table Construction}

\subsection{Isolated DOPC Conformer Ensemble}

The production CGL-CGL and CGL-particle tables are generated from an isolated
DOPC conformer ensemble rather than from a membrane trajectory.  Starting from
the dry-MARTINI DOPC topology, bead coordinates are sampled by Metropolis
moves under the dry-MARTINI bonded energy only.  The current installed
\texttt{dopc.h5} uses two isolated bonded conformers sampled at the converted
experimental production temperature,
\begin{equation}
  T_{\mathrm{prod}}=0.8647~T_{\mathrm{up}} \simeq 303~\mathrm{K}.
\end{equation}
The table metadata records the isolated-conformer provenance explicitly:
\begin{quote}
\small
\texttt{dopc\_reference\_source=isolated\_dopc\_itp\_bonded\_mc\_ensemble}\\
\texttt{dopc\_reference\_conformer\_count=2}.
\end{quote}

This choice is important for transferability.  The CGL-CGL table contains no
bilayer PMF, bilayer force matching, iterative Boltzmann inversion, density
target, bilayer thickness target, or protein-alignment target.  The bilayer is
therefore a prediction of the local projected dry-MARTINI force field, not a
feature imposed by fitting to membrane observables.

\subsection{Tempered Free-Energy Projection}

For a coarse coordinate $x$, let $U(x,\omega)$ be the explicit dry-MARTINI
bead energy for hidden conformer, rotamer, bead-frame, and azimuthal variables
collectively denoted by $\omega$.  The projected table value is a tempered
free energy,
\begin{equation}
  F_{\tau_{\mathrm{avg}}}(x)=
  -k_B\tau_{\mathrm{avg}}
  \log
  \frac{\sum_\omega w_\omega
  \exp[-U(x,\omega)/(k_B\tau_{\mathrm{avg}})]}
  {\sum_\omega w_\omega},
  \qquad
  \tau_{\mathrm{avg}}=25.0 .
\end{equation}
The same hidden-state reduction is used for CGL-CGL, CGL-SC, CGL-particle, and
SC-particle tables.  The hidden coordinates differ by family, but the
underlying contract is identical: explicit dry-MARTINI constituent projection,
the same Lennard-Jones/Coulomb energy helper, the same cutoff and numerical
guard, and the same native-unit to Upside-unit conversion.

For the CGL tables, a conditioned control variable is stored rather than the
raw free energy.  With $E(x)=F_{\tau_{\mathrm{avg}}}(x)$ and
$E_0=\min_x E(x)$ for one type pair,
\begin{equation}
  y(x)=\log\left[1+\frac{E(x)-E_0}{k_BT_{\mathrm{prod}}}\right],
\end{equation}
and runtime applies the exact inverse,
\begin{equation}
  E(x)=E_0+k_BT_{\mathrm{prod}}\left[\exp(y(x))-1\right].
\end{equation}
This transformation improves numerical conditioning while preserving the
tabulated free-energy surface.

\subsection{Interaction Families}

The installed HDF5 files contain four relevant table families.  Their source
objects differ, but their derivation method is the same explicit
dry-MARTINI constituent projection summarized above.

\begin{table}[h]
\centering
\caption{Production HDF5 table families.  CGL tables are stored in
\texttt{parameters/dryMARTINI/dopc.h5}; the SC-particle table is stored in
\texttt{parameters/dryMARTINI/sidechain.h5}.}
\small
\begin{tabular}{p{0.15\linewidth}p{0.25\linewidth}p{0.24\linewidth}p{0.25\linewidth}}
\toprule
Family & Source object & Target object & Runtime representation \\
\midrule
CGL-CGL &
DOPC CGL bead ensemble &
DOPC CGL bead ensemble &
log1p tensor B-spline \\
CGL-SC &
side-chain rotamer bead ensemble &
DOPC CGL bead ensemble &
log1p tensor B-spline \\
CGL-particle &
DOPC CGL bead ensemble &
dry-MARTINI particle type &
log1p radial-angular B-spline \\
SC-particle &
side-chain rotamer bead ensemble &
dry-MARTINI particle type &
direct rotamer grid \\
\bottomrule
\end{tabular}
\label{tab:h5families}
\end{table}

The distinction in the final column is representational, not a change in the
physical source model.  SC-particle interactions are evaluated on a direct
rotamer free-energy grid because they are used by the existing side-chain
rotamer machinery.  The CGL families use spline tables because the CGL runtime
requires differentiable radial/orientational surfaces for forces and torques.

\section{CGL Transport Model}

\subsection{Generalized Langevin Thermostat}

The conservative HDF5 tables determine the potential energy and torques.  They
do not set the physical time scale of the one-particle CGL dynamics.  CGL
transport is therefore controlled by a particle-level generalized Langevin
thermostat applied to CGL translational velocities.  For the fixed two-mode
kernel used here,
\begin{equation}
  K(t)=\sum_{m=1}^2 g_m \exp(-t/\tau_m),
  \qquad
  \tau_1=0.2,\quad \tau_2=2.0 .
\end{equation}
The random force is paired with the memory kernel by fluctuation--dissipation,
\begin{equation}
  \left\langle \eta_\alpha(t)\eta_\beta(t') \right\rangle
  = k_BT_{\mathrm{prod}}K(|t-t'|)\delta_{\alpha\beta}.
\end{equation}
The GLE couplings are global CGL transport parameters.  They are not
interaction-specific force-field couplings and are not stored separately for
CGL-CGL, CGL-SC, CGL-particle, or SC-particle tables.  In the selected model,
\begin{equation}
  (g_1,g_2)=(0.30375,0.2205),
\end{equation}
with CGL mass scale $0.012$, rotational thermostat time $0.008$, and no
explicit ions.

\subsection{Diffusion Target}

The Upside-core clock maps one integration step to approximately
$40~\mathrm{ps}$.  A DOPC lateral diffusion reference of
$14~\mu\mathrm{m}^2\,\mathrm{s}^{-1}$ is reduced by the 56-fold
coarse-graining factor associated with mapping 14 dry-MARTINI beads to one CGL
particle:
\begin{equation}
  D_{\mathrm{target}} =
  \frac{14}{56}
  =0.25~\mu\mathrm{m}^2\,\mathrm{s}^{-1}.
\end{equation}
This target is used only to choose the CGL transport parameters after the
conservative force field has already been fixed.

\section{Parameter Selection and Validation}

\subsection{CGL-Only Diffusion Scan}

The final parameter choice was made from 288-lipid CGL-only bilayers simulated
at $T_{\mathrm{prod}}=0.8647$ with semi-isotropic NPT enabled.  The selected
coupling scale is $0.75$ relative to the reference kernel
$(0.405,0.294)$, giving $(0.30375,0.2205)$.  This point gives
\begin{equation}
  D_{40}=0.25696~\mu\mathrm{m}^2\,\mathrm{s}^{-1},
\end{equation}
where $D_{40}$ denotes the lateral diffusion coefficient interpreted with the
Upside-core $40~\mathrm{ps}$ per step mapping.

\begin{figure}[h]
\centering
\includegraphics[width=0.78\linewidth]{outputs/cgl_timescale_parameter_report/fig1_final_isolated_mc2_diffusion_vs_gle_scale.png}
\caption{CGL lateral diffusion as a function of the two-mode GLE coupling
scale for the strict isolated-conformer DOPC HDF5 table.  The selected point is
the closest stable CGL-only point to the target diffusion.}
\label{fig:diffusion-scan}
\end{figure}

\begin{table}[h]
\centering
\caption{Final isolated-conformer 100k NPT CGL-only scan.  All points pass the
CGL bilayer geometry gate with no orientation-sign leaflet crossing.}
\begin{tabular}{cccc}
\toprule
GLE scale & Couplings $(g_1,g_2)$ & $D_{40}$ &
minimum $|\hat{\mathbf{n}}\cdot\hat{\mathbf{z}}|$ \\
\midrule
0.70 & (0.2835, 0.2058) & 0.2899 & 0.9607 \\
0.75 & (0.30375, 0.2205) & 0.2570 & 0.9605 \\
0.80 & (0.3240, 0.2352) & 0.2400 & 0.9600 \\
0.90 & (0.3645, 0.2646) & 0.1756 & 0.9613 \\
1.00 & (0.4050, 0.2940) & 0.1918 & 0.9626 \\
1.10 & (0.4455, 0.3234) & 0.1526 & 0.9591 \\
1.20 & (0.4860, 0.3528) & 0.1263 & 0.9574 \\
\bottomrule
\end{tabular}
\label{tab:diffusion-scan}
\end{table}

\subsection{Stokes--Einstein Check}

For a fixed-shape GLE kernel, the Stokes--Einstein check should be plotted as
diffusion versus relative mobility, approximated here by the inverse coupling
scale.  The target-window points at scales $0.70$, $0.75$, and $0.80$ are
locally linear with $R^2=0.979$.  The broader single-seed scan has
$R^2=0.947$ and is visibly non-monotonic around scales $0.90$ and $1.00$.
We therefore use the local Stokes--Einstein behavior to support the current
timescale choice, but do not claim robust broad-range Stokes--Einstein
linearity without additional replicate or longer transport simulations.

\begin{figure}[h]
\centering
\includegraphics[width=0.78\linewidth]{outputs/cgl_timescale_parameter_report/fig5_stokes_einstein_local_target_window.png}
\caption{Local Stokes--Einstein check near the selected CGL transport
parameter.  The abscissa is relative mobility, $1/\mathrm{scale}$, rather than
the coupling scale itself.}
\label{fig:se-local}
\end{figure}

\begin{figure}[h]
\centering
\includegraphics[width=0.78\linewidth]{outputs/cgl_timescale_parameter_report/fig6_stokes_einstein_broad_scan.png}
\caption{Broad Stokes--Einstein diagnostic.  The non-monotonic middle points
show that the present single-seed broad scan should not be interpreted as a
publication-quality broad linearity validation.}
\label{fig:se-broad}
\end{figure}

\subsection{Protein--Bilayer Hybrid Tests}

After passing the CGL-only gate, the same HDF5 tables and transport parameters
were used without retuning in 1RKL and 1AFO hybrid protein--bilayer
simulations.  Both direct runs and independent-seed replicates completed stage
7 production with exact $T=0.8647$, finite energies, active CGL-CGL, CGL-SC,
and CGL-particle spline nodes, stable CGL orientations, and no
orientation-sign leaflet crossings.  In the independent replicates, 1RKL
changed from 27.70 to 31.05 total hydrogen-bond score with protein radius of
gyration from 12.68 to $13.12~\text{\AA}$, while 1AFO changed from 84.86 to
82.35 with radius of gyration from 15.82 to $15.37~\text{\AA}$.  The minimum
same-leaflet CGL spacing remained above $6.25~\text{\AA}$ in both replicated
hybrid systems.

\section{Model Scope}

The model contains no twisting coordinate, no hidden orientation marker
particle, no arbitrary force cap, no bilayer-derived production CGL-CGL
correction, and no added orientation restraint.  Protein side-chain and
backbone interactions with the dry-MARTINI environment remain active through
their HDF5 spline tables and backbone proxy projection.  All conservative CGL
interactions used during simulation are table-driven.

The current model should therefore be viewed as a physically constrained
one-particle DOPC representation with a calibrated transport kernel.  The
timescale match is satisfactory at the selected target point, and the
near-target Stokes--Einstein relation is acceptable.  A stronger claim of
broad Stokes--Einstein linearity will require replicate or longer CGL-only
transport simulations over the wider mobility range.

\begin{thebibliography}{9}

\bibitem{MARTINI}
Marrink, S.~J., Risselada, H.~J., Yefimov, S., Tieleman, D.~P., and de Vries, A.~H.
``The MARTINI Force Field: Coarse Grained Model for Biomolecular Simulations.''
\textit{J. Phys. Chem. B}, 111(27), 7812--7824 (2007).

\bibitem{KirkwoodPMF}
Kirkwood, J.~G.
``Statistical Mechanics of Fluid Mixtures.''
\textit{J. Chem. Phys.}, 3, 300--313 (1935).

\bibitem{deBoor}
de Boor, C.
``A Practical Guide to Splines.''
Springer (1978).

\bibitem{Upside}
Jumper, J.~M., Faruk, N.~F., Freed, K.~F., and Sosnick, T.~R.
``Accurate calculation of side chain packing and free energy with
applications to protein molecular dynamics.''
\textit{PLoS Comput. Biol.}, 14(12), e1006342 (2018).

\end{thebibliography}

\end{document}
